Converter-interfaced wind turbines add no natural inertia, and grid-following (GFL) converters that synchronise through a phase-locked loop lose stability as the short-circuit ratio (SCR) falls. This report presents Step 1 of a project on a grid-forming (GFM) PMSG wind turbine that holds a 10 percent power reserve by running on the overspeed side of the $C_p$-optimal line, so that rotor kinetic energy and reserve are released together. We built an aerodynamic model ($C_{p,max}=0.480$ at $\lambda_{opt}=8.10$, below the Betz limit 0.593), a frequency-dynamics model with 40 percent wind, and a quasi-static weak-grid eigenvalue study; all 16 verification gates pass. After a 50 MW load step the GFM plant lifts the nadir from 49.659 Hz (GFL, MPPT) to 49.906 Hz, inside the IEGC band, halves the instantaneous RoCoF (0.162 against 0.332 Hz/s) and raises system inertia from 3.76 s to 8.51 s. The reserve, not the control form, sets the nadir: GFL with the same reserve reaches the same value. The decisive GFM advantage is weak-grid stability: GFL loses stability at SCR 1.30 to 1.37, GFM stays stable to the static limit near 1.0. Step 2 adds detailed dq converter models, current limiting and realistic frequency-measurement windows.
